% 课程要求：单栏、五号字（10.5pt），行距与 Word 模板接近
\usepackage[left=2.54cm,right=2.54cm,top=2.54cm,bottom=2.54cm]{geometry}
\usepackage{setspace}
\setstretch{1.25}

\usepackage{amsmath,amssymb,amsfonts,amsthm}
\usepackage{mathtools}
\usepackage{bm}
\usepackage{graphicx}
\usepackage{booktabs}
\usepackage{tabularx}
\usepackage{longtable}
\usepackage{multirow}
\usepackage{float}
\usepackage{caption}
\usepackage{subcaption}
\usepackage{enumitem}
\usepackage{algorithm}
\usepackage{algpseudocode}
\usepackage[numbers,sort&compress]{natbib}
\usepackage{hyperref}
\usepackage{xcolor}

\hypersetup{
  colorlinks=true,
  linkcolor=black,
  citecolor=blue!60!black,
  urlcolor=blue!60!black
}

\graphicspath{{../figures/}}

% 定理环境（中文）
\theoremstyle{plain}
\newtheorem{theorem}{定理}[section]
\newtheorem{lemma}[theorem]{引理}
\newtheorem{proposition}[theorem]{命题}
\theoremstyle{definition}
\newtheorem{example}[theorem]{示例}
\theoremstyle{remark}
\newtheorem{remark}[theorem]{注}

% 算法标题
\floatname{algorithm}{算法}
\renewcommand{\algorithmicrequire}{\textbf{输入:}}
\renewcommand{\algorithmicensure}{\textbf{输出:}}

% 图表标题
\captionsetup{font=small,labelfont=bf}

% 节标题格式
\ctexset{
  section = {format=\large\bfseries},
  subsection = {format=\normalsize\bfseries},
  subsubsection = {format=\normalsize\bfseries}
}

% 公式编号按节
\numberwithin{equation}{section}

% 页眉页脚
\usepackage{fancyhdr}
\pagestyle{fancy}
\fancyhf{}
\fancyhead[C]{\small 数值分析与算法 · 期末报告}
\fancyfoot[C]{\thepage}
\renewcommand{\headrulewidth}{0.4pt}
